\begin{lemma}[Failure and legal moves in the copied triangle]
Let $h$ be bad for $\mathsf{PowerRel}(r)$.  For every increasing
enumeration $x$, row $n$, and depth $k$, the copied pair is failed and the
next left value is a legal move:
\[
\begin{split}
&\neg\mathsf{PowerRel}(r)(L_{n,k},L_{n+1,k}),\\
&\mathsf{PowerMove}(L_{n,k},L_{n,k+1}),
\end{split}
\]
where $L_{n,k}=\mathsf{PowerCopiedLeft}_r(h,x,n,k)$.
\end{lemma}
\begin{proof}
First consider depth zero.  The pair of initial values is
\bigl(h(\mathsf{PowerShiftIter}(n,x)),
h(\mathsf{PowerShiftIter}(n+1,x))\bigr).  Repeatedly applying the shift
identity from \noderef{PowerShiftIter} identifies the second value
with the shift of the first, so its failure follows from
\noderef{BadMultiSequence}.  The canonical response is legal by
\noderef{PowerIResponseLegal}.

Assume the two assertions at depth $k$ for every row.  The second assertion
at row $n+1$ says that $L_{n+1,k+1}$ is a legal response to
$L_{n+1,k}$.  Apply \noderef{PowerIResponseFailure} to the failed pair
$(L_{n,k},L_{n+1,k})$; it gives failure of
$(L_{n,k+1},L_{n+1,k+1})$.  Applying
\noderef{PowerIResponseLegal} to that failed pair gives the next legal left
move.  Induction on $k$ proves both claims simultaneously for all rows.
\end{proof}
